\section{Curvature bias of greyscale marching cubes}\label{app:curvature}
For a plane and a symmetric blur the
$0.5$ iso-surface of the blurred image coincides with the true surface. For a curved surface it does
not. If the effective blur is Gaussian with variance $\sigma_\text{eff}^2$, the level sets of the
blurred indicator move towards the centre of curvature by
\begin{equation}
  \delta \approx \sigma_\text{eff}^2\,H, \qquad H = \tfrac12(\kappa_1+\kappa_2),
  \label{eq:shift}
\end{equation}
where $H$ is the mean curvature. The intuition is that blurring a curved object spreads its intensity
outwards over a larger volume, so the half-intensity surface ends up slightly inside the true one, and
more so where the surface bends sharply. (Formally this follows from the heat equation, since Gaussian
blurring is diffusion for a time $\sigma_\text{eff}^2/2$.) For a sphere of radius $R$ this gives a relative area error of
$-2\delta/R = -2\sigma_\text{eff}^2/R^2$. In our forward model the voxel box ($\sigma^2 = 1/12$) and the
Gaussian PSF ($\sigma = 0.6$) combine to $\sigma_\text{eff}^2 \approx 0.443$ voxel$^2$; adding a
smoothing filter with $\sigma_s = 1$ raises it to $1.443$. Table~\ref{tab:curvature} compares this
prediction with the measured errors on clean spheres. The agreement is close at low resolution, where
the curvature term dominates; at $64^3$ the remaining discrepancy is of the same size as the
interpolation error of marching cubes itself.

\begin{table}[h]
  \centering\footnotesize\setlength{\tabcolsep}{3pt}
  \caption{Greyscale marching-cubes error on clean spheres: measured vs the curvature prediction
  $-2\sigma_\text{eff}^2/R^2$ from Eq.~\eqref{eq:shift}. $R$ in voxels.}
  \label{tab:curvature}
  \begin{tabular}{rrrrrr}
    \toprule
    & & \multicolumn{2}{c}{MC on greyscale} & \multicolumn{2}{c}{MC after smoothing ($\sigma_s=1$)}\\
    \cmidrule(lr){3-4}\cmidrule(lr){5-6}
    Grid & $R$ & predicted & measured & predicted & measured \\
    \midrule
    $16^3$ & 5.4  & $-3.04\%$ & $-3.65\%$ & $-9.90\%$ & $-10.99\%$ \\
    $32^3$ & 12.6 & $-0.56\%$ & $-0.62\%$ & $-1.82\%$ & $-2.00\%$ \\
    $64^3$ & 27.0 & $-0.12\%$ & $-0.05\%$ & $-0.40\%$ & $-0.43\%$ \\
    $128^3$ & 55.8 & $-0.03\%$ & $+0.07\%$ & $-0.09\%$ & $-0.10\%$ \\
    \bottomrule
  \end{tabular}
\end{table}

\section{Reference surface areas}\label{app:reference}
Spheres, cylinders and tori have closed-form areas. For an ellipsoid with semi-axes $a\ge b\ge c$,
\begin{equation}
\begin{aligned}
  A &= 2\pi c^2 + \frac{2\pi ab}{\sin\varphi}\left[E(\varphi,m)\sin^2\varphi + F(\varphi,m)\cos^2\varphi\right],\\
  \cos\varphi &= \frac ca,\qquad m = \frac{a^2(b^2-c^2)}{b^2(a^2-c^2)},
\end{aligned}
\end{equation}
with $F$ and $E$ the incomplete elliptic integrals of the first and second kind. The remaining
families (random blobs and superellipsoids) are star-shaped, i.e.\ their surface is
$\{r(\mathbf u)\,\mathbf u : |\mathbf u|=1\}$ for a radial function $r$. We triangulate them by mapping
the vertices of a subdivided icosahedron to radius $r(\mathbf u)$. For a smooth surface the area of such
an inscribed mesh converges as $A_\ell = A + c\,h_\ell^2 + o(h_\ell^2)$ in the edge length $h_\ell$, and
each subdivision halves $h_\ell$, so Richardson extrapolation of two levels,
\begin{equation}
  A_\text{ref} = \tfrac13\left(4A_{\ell+1} - A_\ell\right),
\end{equation}
cancels the leading error term. We use levels 6 and 7 (81,920 and 327,680 triangles) and keep
$|A_\text{ref}-A_7|$ as an error estimate for each shape.
